\chapter{Requerimientos de Usuario}

% ============================================================
\section{Diagrama de Casos de Uso}

Para mayor claridad, el diagrama de casos de uso se divide en tres subsistemas.

% --- Subsistema 1: Solicitante ---
\subsection{Subsistema: Solicitante (Alumno / Docente)}

\begin{figure}[H]
\centering
\begin{tikzpicture}[scale=0.9, transform shape,
    uc/.style={draw, ellipse, align=center, font=\footnotesize, inner sep=3pt, minimum width=3cm, minimum height=1cm},
    act/.style={font=\small\bfseries, align=center},
    sbound/.style={draw, dashed, rounded corners, inner sep=0.8cm},
    inc/.style={dashed, -{Stealth}, font=\scriptsize},
    generalization/.style={-{Triangle[open, length=3mm, width=2.5mm]}, thick}
]

% Actor abstracto (Solicitante) centrado verticalmente
\node[act] (solicitante) at (-3.5, -5) {Solicitante};

% Actores concretos
\node[act] (alumno) at (-6.5, -2.5) {Alumno};
\node[act] (docente) at (-6.5, -7.5) {Docente};

% Generalizacion: Alumno y Docente -> Solicitante
\draw[generalization] (alumno) -- (solicitante);
\draw[generalization] (docente) -- (solicitante);

% Casos de uso (espaciado de 2.5 unidades)
\node[uc] (cu01) at (2.5, 0) {CU-01\\Crear solicitud};
\node[uc] (cu02) at (2.5, -2.5) {CU-02\\Consultar solicitudes};
\node[uc] (cu03) at (2.5, -5) {CU-03\\Ver detalle};
\node[uc] (cu04) at (2.5, -7.5) {CU-04\\Cancelar solicitud};
\node[uc] (cu05) at (2.5, -10) {CU-05\\Ver seguimiento};

% CU-17 a la derecha
\node[uc] (cu17) at (7.5, 0) {CU-17\\Notificar por correo};

% Limites del sistema
\node[sbound, fit=(cu01)(cu02)(cu03)(cu04)(cu05)(cu17), label={[font=\bfseries\small]above:Sistema de Solicitudes -- Solicitante}] (boundary) {};

% Conexiones Solicitante -> todos los CU
\draw[thick] (solicitante) -- (cu01.west);
\draw[thick] (solicitante) -- (cu02.west);
\draw[thick] (solicitante) -- (cu03.west);
\draw[thick] (solicitante) -- (cu04.west);
\draw[thick] (solicitante) -- (cu05.west);

% Include CU-01 -> CU-17
\draw[inc] (cu01) -- node[above, font=\scriptsize] {$\ll$include$\gg$} (cu17);

\end{tikzpicture}
\caption{Diagrama de casos de uso -- Subsistema Solicitante.}
\label{fig:uc-solicitante}
\end{figure}

% --- Subsistema 2: Atención ---
\subsection{Subsistema: Atención de Solicitudes (Personal)}

\begin{figure}[H]
\centering
\begin{tikzpicture}[scale=0.9, transform shape,
    uc/.style={draw, ellipse, align=center, font=\footnotesize, inner sep=3pt, minimum width=3cm, minimum height=1cm},
    act/.style={font=\small\bfseries, align=center},
    sbound/.style={draw, dashed, rounded corners, inner sep=0.8cm},
    inc/.style={dashed, -{Stealth}, font=\scriptsize},
    ext/.style={dashed, -{Stealth}, font=\scriptsize},
    generalization/.style={-{Triangle[open, length=3mm, width=2.5mm]}, thick}
]

% Actor abstracto (Personal) centrado
\node[act] (personal) at (-3.5, -5) {Personal};

% Actores concretos
\node[act] (ce) at (-6.5, -2.5) {Control\\Escolar};
\node[act] (rp) at (-6.5, -7.5) {Responsable\\de Programa};

% Generalizacion
\draw[generalization] (ce) -- (personal);
\draw[generalization] (rp) -- (personal);

% Casos de uso principales (columna izquierda, espaciado 2.5)
\node[uc] (cu06) at (2.5, 0) {CU-06\\Listar asignadas};
\node[uc] (cu07) at (2.5, -2.5) {CU-07\\Atender solicitud};
\node[uc] (cu08) at (2.5, -5) {CU-08\\Cambiar estado};
\node[uc] (cu03b) at (2.5, -7.5) {CU-03\\Ver detalle};

% Casos de uso secundarios (columna derecha, alíneados)
\node[uc] (cu09) at (7.5, -2.5) {CU-09\\Generar PDF};
\node[uc] (cu17b) at (7.5, -5) {CU-17\\Notificar por correo};

% Limites
\node[sbound, fit=(cu06)(cu07)(cu08)(cu09)(cu03b)(cu17b), label={[font=\bfseries\small]above:Sistema de Solicitudes -- Atención}] (boundary) {};

% Conexiones Personal -> CU
\draw[thick] (personal) -- (cu06.west);
\draw[thick] (personal) -- (cu07.west);
\draw[thick] (personal) -- (cu08.west);
\draw[thick] (personal) -- (cu03b.west);

% Extend CU-09 -> CU-07
\draw[ext] (cu09) -- node[above, font=\scriptsize] {$\ll$extend$\gg$} (cu07);

% Include CU-08 -> CU-17
\draw[inc] (cu08) -- node[above, font=\scriptsize] {$\ll$include$\gg$} (cu17b);

\end{tikzpicture}
\caption{Diagrama de casos de uso -- Subsistema Atención.}
\label{fig:uc-atención}
\end{figure}

% --- Subsistema 3: Administración ---
\subsection{Subsistema: Administración}

\begin{figure}[H]
\centering
\begin{tikzpicture}[scale=0.9, transform shape,
    uc/.style={draw, ellipse, align=center, font=\footnotesize, inner sep=3pt, minimum width=3cm, minimum height=1cm},
    act/.style={font=\small\bfseries, align=center},
    sbound/.style={draw, dashed, rounded corners, inner sep=0.8cm},
    inc/.style={dashed, -{Stealth}, font=\scriptsize}
]

% Actor
\node[act] (admin) at (-3.5, -3.75) {Administrador};

% Casos de uso (columna izquierda, espaciado 2.5)
\node[uc] (cu10) at (2.5, 0) {CU-10\\Gestiónar tipos};
\node[uc] (cu11) at (2.5, -2.5) {CU-11\\Gestiónar formularios};
\node[uc] (cu12) at (2.5, -5) {CU-12\\Gestiónar plantillas};

% Casos de uso (columna derecha)
\node[uc] (cu13) at (7.5, 0) {CU-13\\Gestiónar mentores};
\node[uc] (cu14) at (7.5, -2.5) {CU-14\\Ver dashboard};
\node[uc] (cu15) at (7.5, -5) {CU-15\\Exportar reportes};

% Limites
\node[sbound, fit=(cu10)(cu11)(cu12)(cu13)(cu14)(cu15), label={[font=\bfseries\small]above:Sistema de Solicitudes -- Administración}] (boundary) {};

% Conexiones
\draw[thick] (admin) -- (cu10.west);
\draw[thick] (admin) -- (cu11.west);
\draw[thick] (admin) -- (cu12.west);
\draw[thick] (admin) -- (cu13.west);
\draw[thick] (admin) -- (cu14.west);
\draw[thick] (admin) -- (cu15.west);

\end{tikzpicture}
\caption{Diagrama de casos de uso -- Subsistema Administración.}
\label{fig:uc-admin}
\end{figure}

% ============================================================
\section{Catálogo de Casos de Uso}

\begin{longtable}{|C{1.2cm}|L{12.5cm}|}
    \hline
    \textbf{ID} & \textbf{Descripción} \\
    \hline
    \endfirsthead
    \hline
    \textbf{ID} & \textbf{Descripción} \\
    \hline
    \endhead
    \hline
    \endfoot

    CU-01 & \textbf{Crear solicitud:} El solicitante (alumno o docente) seleccióna un tipo de solicitud, llena el formulario dinámico, adjunta archivos requeridos y envia la solicitud. \\
    \hline
    CU-02 & \textbf{Consultar solicitudes propias:} El solicitante visualiza un listado de todas las solicitudes que ha creado, con su folio, tipo y estado actual. \\
    \hline
    CU-03 & \textbf{Ver detalle de solicitud:} Cualquier usuario autorizado puede ver la información completa de una solicitud: datos del formulario, archivos adjuntos e historial de seguimiento. \\
    \hline
    CU-04 & \textbf{Cancelar solicitud propia:} El solicitante cancela una solicitud que se encuentre en estado ``Creada''. \\
    \hline
    CU-05 & \textbf{Ver seguimiento de solicitud:} El solicitante consulta el historial de cambios de estado de su solicitud con fechas y observaciones. \\
    \hline
    CU-06 & \textbf{Listar solicitudes asignadas:} El personal visualiza las solicitudes dirigidas a su rol, con filtros por estado, tipo, fecha y busqueda por folio o nombre. \\
    \hline
    CU-07 & \textbf{Atender solicitud:} El personal revisa el detalle de una solicitud asignada a su rol y realiza las acciones necesarias para su resolución. \\
    \hline
    CU-08 & \textbf{Cambiar estado de solicitud:} El personal modifica el estado de una solicitud (En proceso, Finalizada o Cancelada) y agrega observaciones. \\
    \hline
    CU-09 & \textbf{Generar documento PDF:} El personal genera un documento PDF a partir de la plantilla asociada al tipo de solicitud, sustituyendo las variables con datos del solicitante. \\
    \hline
    CU-10 & \textbf{Gestiónar tipos de solicitud:} El administrador crea, edita o elimina tipos de solicitud del catálogo. \\
    \hline
    CU-11 & \textbf{Gestiónar formularios dinámicos:} El administrador configura los campos del formulario asociado a cada tipo de solicitud. \\
    \hline
    CU-12 & \textbf{Gestiónar plantillas:} El administrador crea, edita o elimina plantillas de documentos con variables sustituibles para tipos de solicitud que lo requieran. \\
    \hline
    CU-13 & \textbf{Gestiónar lista de mentores:} El administrador agrega, elimina o consulta alumnos mentores activos de forma manual o mediante carga de archivo CSV. \\
    \hline
    CU-14 & \textbf{Ver dashboard y métricas:} El administrador visualiza estadísticas de solicitudes por tipo, estado y período. \\
    \hline
    CU-15 & \textbf{Exportar reportes:} El administrador exporta las métricas en formato CSV o PDF. \\
    \hline
    CU-17 & \textbf{Enviar notificación por correo:} El sistema envia automáticamente correos electrónicos al personal cuando se crea una solicitud y al solicitante cuando cambia el estado. \\
    \hline
\end{longtable}

% ============================================================
\section{Especificación de Casos de Uso}

% --- CU-01 ---
\subsection{CU-01: Crear solicitud}

\begin{longtable}{|L{4cm}|L{9.5cm}|}
    \hline
    \multicolumn{2}{|c|}{\textbf{CU-01: Crear solicitud}} \\
    \hline
    \endfirsthead
    \hline
    \endfoot

    \textbf{Fecha y Versión:} & Marzo 2026, v1.0 \\
    \hline
    \textbf{Autor:} & Equipo de desarrollo \\
    \hline
    \textbf{Descripción:} & El solicitante (alumno o docente) crea una nueva solicitud selecciónando un tipo del catálogo, llenando el formulario dinámico correspondiente y adjuntando los archivos requeridos. El sistema genera un folio único, registra la solicitud en estado ``Creada'' y notifica al personal responsable por correo electrónico. \\
    \hline
    \textbf{Actor(es) primario(s):} & Alumno, Docente \\
    \hline
    \textbf{Actor(es) secundario(s):} & Proveedor de autenticación (JWT), SIGA (datos del solicitante), Servicio SMTP (notificación) \\
    \hline
    \textbf{Precondición(es):} &
    \begin{enumerate}[nosep, leftmargin=*]
        \item El usuario esta autenticado mediante token JWT válido.
        \item Existe al menos un tipo de solicitud configurado en el catálogo.
        \item El tipo de solicitud tiene un formulario dinámico asociado con al menos un campo.
    \end{enumerate} \\
    \hline
    \textbf{Happy Path:} &
    \begin{enumerate}[nosep, leftmargin=*]
        \item El solicitante accede a la sección ``Crear solicitud''.
        \item El sistema muestra la lista de tipos de solicitud disponibles para el rol del usuario.
        \item El solicitante seleccióna un tipo de solicitud.
        \item El sistema presenta el formulario dinámico correspondiente al tipo selecciónado con todos sus campos configurados (texto, fecha, archivo, etc.).
        \item El solicitante llena los campos del formulario y adjunta los archivos requeridos.
        \item El solicitante presiona el boton ``Enviar solicitud''.
        \item El sistema valida que todos los campos requeridos esten completos y que los archivos cumplan con los formatos y tamaños permitidos.
        \item El sistema genera un folio único para la solicitud.
        \item El sistema registra la solicitud en estado ``Creada'' con la fecha actual.
        \item El sistema crea el registro inicial de seguimiento.
        \item El sistema envia una notificación por correo al personal con el rol responsable del tipo de solicitud (CU-17).
        \item El sistema muestra al solicitante un mensaje de confirmación con el folio generado.
    \end{enumerate} \\
    \hline
    \textbf{Postcondición(es):} &
    \begin{enumerate}[nosep, leftmargin=*]
        \item La solicitud queda registrada en el sistema con estado ``Creada'' y folio único.
        \item Los archivos adjuntos quedan almacenados y asociados a la solicitud.
        \item El personal responsable recibe una notificación por correo.
        \item El registro de seguimiento inicial esta creado.
    \end{enumerate} \\
    \hline
    \textbf{Alternativos:} &
    \textbf{FA-01: El tipo de solicitud requiere pago.}
    \begin{enumerate}[nosep, leftmargin=*]
        \item[4.1] El sistema detecta que el tipo requiere comprobante de pago.
        \item[4.2] El sistema verifica si el solicitante es un alumno mentor activo.
        \item[4.3a] Si es mentor: el campo de comprobante de pago se omite o se marca como no requerido. Continua en paso 5.
        \item[4.3b] Si no es mentor: el formulario incluye un campo obligatorio para adjuntar comprobante de pago. Continua en paso 5.
    \end{enumerate}
    \vspace{0.2cm}
    \textbf{FA-02: Validación falla (campos incompletos).}
    \begin{enumerate}[nosep, leftmargin=*]
        \item[7.1] El sistema detecta que hay campos requeridos sin completar o archivos faltantes.
        \item[7.2] El sistema muestra mensajes de error indicando los campos que deben corregirse.
        \item[7.3] El solicitante corrige la información.
        \item[7.4] Regresa al paso 6.
    \end{enumerate}
    \vspace{0.2cm}
    \textbf{FA-03: El archivo excede el tamaño permitido.}
    \begin{enumerate}[nosep, leftmargin=*]
        \item[7.1] El sistema detecta que un archivo excede el tamaño máximo.
        \item[7.2] El sistema muestra un mensaje indicando el límite de tamaño.
        \item[7.3] El solicitante adjunta un archivo más pequeno.
        \item[7.4] Regresa al paso 6.
    \end{enumerate}
    \vspace{0.2cm}
    \textbf{FA-04: Error al enviar notificación.}
    \begin{enumerate}[nosep, leftmargin=*]
        \item[11.1] El servicio SMTP no esta disponible.
        \item[11.2] El sistema registra el error internamente.
        \item[11.3] La solicitud se crea éxitosamente, pero la notificación queda pendiente.
        \item[11.4] Continua en paso 12.
    \end{enumerate} \\
    \hline
    \textbf{Notas:} & Los tipos de solicitud disponibles dependen del rol del usuario (alumno o docente). Los archivos se almacenan organizados por folio. \\
    \hline
    \textbf{Req. relacionados:} & RF-01, RF-02, RF-03, RF-04, RF-07, RF-10, RF-11 \\
    \hline
\end{longtable}

% --- CU-02 ---
\subsection{CU-02: Consultar solicitudes propias}

\begin{longtable}{|L{4cm}|L{9.5cm}|}
    \hline
    \multicolumn{2}{|c|}{\textbf{CU-02: Consultar solicitudes propias}} \\
    \hline
    \endfirsthead
    \hline
    \endfoot

    \textbf{Fecha y Versión:} & Marzo 2026, v1.0 \\
    \hline
    \textbf{Autor:} & Equipo de desarrollo \\
    \hline
    \textbf{Descripción:} & El solicitante consulta la lista de todas las solicitudes que ha creado, visualizando su folio, tipo, fecha de creación y estado actual. Puede filtrar por estado. \\
    \hline
    \textbf{Actor(es) primario(s):} & Alumno, Docente \\
    \hline
    \textbf{Actor(es) secundario(s):} & --- \\
    \hline
    \textbf{Precondición(es):} &
    \begin{enumerate}[nosep, leftmargin=*]
        \item El usuario esta autenticado.
        \item El usuario tiene al menos una solicitud creada.
    \end{enumerate} \\
    \hline
    \textbf{Happy Path:} &
    \begin{enumerate}[nosep, leftmargin=*]
        \item El solicitante accede a la sección ``Mis solicitudes''.
        \item El sistema muestra una lista páginada de las solicitudes del usuario, ordenadas por fecha de creación (más recientes primero).
        \item Cada elemento de la lista muestra: folio, tipo de solicitud, fecha de creación y estado actual (con indicador visual de color).
        \item El solicitante puede filtrar la lista por estado (Creada, En proceso, Finalizada, Cancelada).
        \item El solicitante seleccióna una solicitud para ver su detalle (CU-03).
    \end{enumerate} \\
    \hline
    \textbf{Postcondición(es):} & El solicitante visualiza la lista de sus solicitudes. \\
    \hline
    \textbf{Alternativos:} &
    \textbf{FA-01: No tiene solicitudes.}
    \begin{enumerate}[nosep, leftmargin=*]
        \item[2.1] El sistema muestra un mensaje indicando que no tiene solicitudes registradas.
        \item[2.2] Se ofrece un enlace para crear una nueva solicitud.
    \end{enumerate} \\
    \hline
    \textbf{Req. relacionados:} & RF-08 \\
    \hline
\end{longtable}

% --- CU-03 ---
\subsection{CU-03: Ver detalle de solicitud}

\begin{longtable}{|L{4cm}|L{9.5cm}|}
    \hline
    \multicolumn{2}{|c|}{\textbf{CU-03: Ver detalle de solicitud}} \\
    \hline
    \endfirsthead
    \hline
    \endfoot

    \textbf{Fecha y Versión:} & Marzo 2026, v1.0 \\
    \hline
    \textbf{Autor:} & Equipo de desarrollo \\
    \hline
    \textbf{Descripción:} & Un usuario autorizado visualiza la información completa de una solicitud: datos del solicitante, respuestas del formulario, archivos adjuntos e historial de seguimiento. \\
    \hline
    \textbf{Actor(es) primario(s):} & Alumno, Docente, Control Escolar, Responsable de Programa, Administrador \\
    \hline
    \textbf{Actor(es) secundario(s):} & SIGA (para obtener datos completos del solicitante) \\
    \hline
    \textbf{Precondición(es):} &
    \begin{enumerate}[nosep, leftmargin=*]
        \item El usuario esta autenticado.
        \item La solicitud existe en el sistema.
        \item El usuario tiene permiso para verla (es el solicitante, tiene el rol responsable, o es administrador).
    \end{enumerate} \\
    \hline
    \textbf{Happy Path:} &
    \begin{enumerate}[nosep, leftmargin=*]
        \item El usuario seleccióna una solicitud de la lista (CU-02 o CU-06).
        \item El sistema verifica que el usuario tiene permiso para ver la solicitud.
        \item El sistema obtiene los datos completos del solicitante desde el SIGA (nombre, programa, etc.).
        \item El sistema muestra la información organizada en secciónes:
            \begin{itemize}[nosep]
                \item Datos generales: folio, tipo, estado actual, fecha de creación.
                \item Datos del solicitante: nombre, matrícula, programa, correo.
                \item Respuestas del formulario: cada campo con su etiqueta y valor.
                \item Archivos adjuntos: lista de archivos con opcion de descarga.
                \item Historial de seguimiento: cambios de estado con fecha, responsable y observaciones.
            \end{itemize}
    \end{enumerate} \\
    \hline
    \textbf{Postcondición(es):} & El usuario visualiza toda la información de la solicitud. \\
    \hline
    \textbf{Alternativos:} &
    \textbf{FA-01: SIGA no disponible.}
    \begin{enumerate}[nosep, leftmargin=*]
        \item[3.1] El sistema no puede conectarse al SIGA.
        \item[3.2] Se muestra la información básica del solicitante (matrícula y correo del token) sin datos adicionales.
    \end{enumerate}
    \vspace{0.2cm}
    \textbf{FA-02: Sin permiso.}
    \begin{enumerate}[nosep, leftmargin=*]
        \item[2.1] El usuario no tiene permiso para ver la solicitud.
        \item[2.2] El sistema muestra un mensaje de acceso denegado.
    \end{enumerate} \\
    \hline
    \textbf{Req. relacionados:} & RF-08, RF-09, RF-10, RNF-02 \\
    \hline
\end{longtable}

% --- CU-04 ---
\subsection{CU-04: Cancelar solicitud propia}

\begin{longtable}{|L{4cm}|L{9.5cm}|}
    \hline
    \multicolumn{2}{|c|}{\textbf{CU-04: Cancelar solicitud propia}} \\
    \hline
    \endfirsthead
    \hline
    \endfoot

    \textbf{Fecha y Versión:} & Marzo 2026, v1.0 \\
    \hline
    \textbf{Autor:} & Equipo de desarrollo \\
    \hline
    \textbf{Descripción:} & El solicitante cancela una solicitud que creo previamente, siempre que esta se encuentre en estado ``Creada''. \\
    \hline
    \textbf{Actor(es) primario(s):} & Alumno, Docente \\
    \hline
    \textbf{Actor(es) secundario(s):} & --- \\
    \hline
    \textbf{Precondición(es):} &
    \begin{enumerate}[nosep, leftmargin=*]
        \item El usuario esta autenticado.
        \item La solicitud pertenece al usuario.
        \item La solicitud esta en estado ``Creada''.
    \end{enumerate} \\
    \hline
    \textbf{Happy Path:} &
    \begin{enumerate}[nosep, leftmargin=*]
        \item El solicitante accede al detalle de una de sus solicitudes (CU-03).
        \item El sistema muestra el boton ``Cancelar solicitud'' (visible solo si el estado es ``Creada'').
        \item El solicitante presiona ``Cancelar solicitud''.
        \item El sistema muestra un dialogo de confirmación: ``¿Esta seguro de que desea cancelar esta solicitud?''.
        \item El solicitante confirma la cancelación.
        \item El sistema cambia el estado de la solicitud a ``Cancelada''.
        \item El sistema registra el cambio en el historial de seguimiento con la fecha, el usuario que cancelo y la observacion ``Cancelada por el solicitante''.
        \item El sistema muestra un mensaje de confirmación.
    \end{enumerate} \\
    \hline
    \textbf{Postcondición(es):} & La solicitud queda en estado ``Cancelada'' con el registro en el historial de seguimiento. \\
    \hline
    \textbf{Alternativos:} &
    \textbf{FA-01: La solicitud no esta en estado ``Creada''.}
    \begin{enumerate}[nosep, leftmargin=*]
        \item[2.1] El boton ``Cancelar solicitud'' no se muestra porque el estado es ``En proceso'', ``Finalizada'' o ``Cancelada''.
        \item[2.2] El solicitante no puede cancelar la solicitud.
    \end{enumerate}
    \vspace{0.2cm}
    \textbf{FA-02: El solicitante no confirma.}
    \begin{enumerate}[nosep, leftmargin=*]
        \item[5.1] El solicitante presiona ``No'' en el dialogo de confirmación.
        \item[5.2] El sistema cierra el dialogo. La solicitud permanece sin cambios.
    \end{enumerate} \\
    \hline
    \textbf{Req. relacionados:} & RF-05 \\
    \hline
\end{longtable}

% --- CU-05 ---
\subsection{CU-05: Ver seguimiento de solicitud}

\begin{longtable}{|L{4cm}|L{9.5cm}|}
    \hline
    \multicolumn{2}{|c|}{\textbf{CU-05: Ver seguimiento de solicitud}} \\
    \hline
    \endfirsthead
    \hline
    \endfoot

    \textbf{Fecha y Versión:} & Marzo 2026, v1.0 \\
    \hline
    \textbf{Autor:} & Equipo de desarrollo \\
    \hline
    \textbf{Descripción:} & El solicitante consulta el historial detallado de cambios de estado de una solicitud propia, incluyendo fechas, quien realizo el cambio y las observaciones asociadas. \\
    \hline
    \textbf{Actor(es) primario(s):} & Alumno, Docente \\
    \hline
    \textbf{Actor(es) secundario(s):} & --- \\
    \hline
    \textbf{Precondición(es):} &
    \begin{enumerate}[nosep, leftmargin=*]
        \item El usuario esta autenticado.
        \item La solicitud pertenece al usuario.
    \end{enumerate} \\
    \hline
    \textbf{Happy Path:} &
    \begin{enumerate}[nosep, leftmargin=*]
        \item El solicitante accede al detalle de una solicitud (CU-03).
        \item El solicitante seleccióna la sección o pestana ``Seguimiento''.
        \item El sistema muestra una línea de tiempo o tabla con el historial:
            \begin{itemize}[nosep]
                \item Fecha y hora del cambio.
                \item Estado anterior y nuevo estado.
                \item Quien realizo el cambio.
                \item Observaciones.
            \end{itemize}
        \item El estado actual se resalta visualmente.
    \end{enumerate} \\
    \hline
    \textbf{Postcondición(es):} & El solicitante visualiza el historial completo de seguimiento. \\
    \hline
    \textbf{Alternativos:} &
    \textbf{FA-01: Solo existe el registro inicial.}
    \begin{enumerate}[nosep, leftmargin=*]
        \item[3.1] Solo se muestra el registro de creación (la solicitud no ha tenido cambios de estado).
    \end{enumerate} \\
    \hline
    \textbf{Req. relacionados:} & RF-05, RF-08 \\
    \hline
\end{longtable}

% --- CU-06 ---
\subsection{CU-06: Listar solicitudes asignadas}

\begin{longtable}{|L{4cm}|L{9.5cm}|}
    \hline
    \multicolumn{2}{|c|}{\textbf{CU-06: Listar solicitudes asignadas}} \\
    \hline
    \endfirsthead
    \hline
    \endfoot

    \textbf{Fecha y Versión:} & Marzo 2026, v1.0 \\
    \hline
    \textbf{Autor:} & Equipo de desarrollo \\
    \hline
    \textbf{Descripción:} & El personal visualiza las solicitudes que han sido dirigidas a su rol, con opciones de filtrado y busqueda para fácilitar la gestión. \\
    \hline
    \textbf{Actor(es) primario(s):} & Control Escolar, Responsable de Programa \\
    \hline
    \textbf{Actor(es) secundario(s):} & --- \\
    \hline
    \textbf{Precondición(es):} &
    \begin{enumerate}[nosep, leftmargin=*]
        \item El usuario esta autenticado con un rol de personal (control escolar o responsable de programa).
    \end{enumerate} \\
    \hline
    \textbf{Happy Path:} &
    \begin{enumerate}[nosep, leftmargin=*]
        \item El personal accede a la sección ``Solicitudes asignadas''.
        \item El sistema consulta las solicitudes cuyo tipo de solicitud tiene como responsable el rol del usuario.
        \item El sistema muestra una lista páginada con: folio, tipo, nombre del solicitante, fecha de creación y estado.
        \item El personal puede filtrar por:
            \begin{itemize}[nosep]
                \item Estado (Creada, En proceso, Finalizada, Cancelada, o Todas).
                \item Tipo de solicitud.
                \item Rango de fechas.
            \end{itemize}
        \item El personal puede buscar por folio o nombre del solicitante.
        \item El sistema muestra contadores de solicitudes por estado.
        \item El personal seleccióna una solicitud para atenderla (CU-07).
    \end{enumerate} \\
    \hline
    \textbf{Postcondición(es):} & El personal visualiza la lista filtrada de solicitudes asignadas a su rol. \\
    \hline
    \textbf{Alternativos:} &
    \textbf{FA-01: No hay solicitudes asignadas.}
    \begin{enumerate}[nosep, leftmargin=*]
        \item[3.1] El sistema muestra un mensaje indicando que no hay solicitudes pendientes.
    \end{enumerate} \\
    \hline
    \textbf{Req. relacionados:} & RF-06, RF-08 \\
    \hline
\end{longtable}

% --- CU-07 ---
\subsection{CU-07: Atender solicitud}

\begin{longtable}{|L{4cm}|L{9.5cm}|}
    \hline
    \multicolumn{2}{|c|}{\textbf{CU-07: Atender solicitud}} \\
    \hline
    \endfirsthead
    \hline
    \endfoot

    \textbf{Fecha y Versión:} & Marzo 2026, v1.0 \\
    \hline
    \textbf{Autor:} & Equipo de desarrollo \\
    \hline
    \textbf{Descripción:} & El personal revisa el detalle completo de una solicitud asignada a su rol, incluyendo los datos del formulario, archivos adjuntos y el historial de seguimiento. A partir de esta revision, decide las acciones a tomar (cambiar estado, generar PDF). \\
    \hline
    \textbf{Actor(es) primario(s):} & Control Escolar, Responsable de Programa \\
    \hline
    \textbf{Actor(es) secundario(s):} & SIGA \\
    \hline
    \textbf{Precondición(es):} &
    \begin{enumerate}[nosep, leftmargin=*]
        \item El usuario esta autenticado con rol de personal.
        \item La solicitud esta asignada al rol del usuario.
    \end{enumerate} \\
    \hline
    \textbf{Happy Path:} &
    \begin{enumerate}[nosep, leftmargin=*]
        \item El personal seleccióna una solicitud de la lista de asignadas (CU-06).
        \item El sistema muestra el detalle completo (CU-03) más opciones de accion.
        \item El personal revisa los datos del formulario y los archivos adjuntos.
        \item El personal descarga los archivos si es necesario.
        \item El personal decide la accion a tomar:
            \begin{itemize}[nosep]
                \item Marcar ``En proceso'' (CU-08).
                \item Finalizar la solicitud (CU-08).
                \item Cancelar la solicitud (CU-08).
                \item Generar documento PDF si aplica (CU-09).
            \end{itemize}
    \end{enumerate} \\
    \hline
    \textbf{Postcondición(es):} & El personal ha revisado la solicitud y puede proceder con las acciones correspondientes. \\
    \hline
    \textbf{Req. relacionados:} & RF-09, RF-10 \\
    \hline
\end{longtable}

% --- CU-08 ---
\subsection{CU-08: Cambiar estado de solicitud}

\begin{longtable}{|L{4cm}|L{9.5cm}|}
    \hline
    \multicolumn{2}{|c|}{\textbf{CU-08: Cambiar estado de solicitud}} \\
    \hline
    \endfirsthead
    \hline
    \endfoot

    \textbf{Fecha y Versión:} & Marzo 2026, v1.0 \\
    \hline
    \textbf{Autor:} & Equipo de desarrollo \\
    \hline
    \textbf{Descripción:} & El personal modifica el estado de una solicitud conforme al ciclo de vida definido y registra observaciones sobre la decision. El sistema notifica al solicitante del cambio. \\
    \hline
    \textbf{Actor(es) primario(s):} & Control Escolar, Responsable de Programa \\
    \hline
    \textbf{Actor(es) secundario(s):} & Servicio SMTP \\
    \hline
    \textbf{Precondición(es):} &
    \begin{enumerate}[nosep, leftmargin=*]
        \item El usuario esta autenticado con rol de personal.
        \item La solicitud esta asignada al rol del usuario.
        \item La solicitud no esta en estado ``Finalizada'' ni ``Cancelada''.
    \end{enumerate} \\
    \hline
    \textbf{Happy Path:} &
    \begin{enumerate}[nosep, leftmargin=*]
        \item El personal visualiza la solicitud (CU-07).
        \item El personal seleccióna la accion de cambio de estado.
        \item El sistema muestra las opciones disponibles segun el estado actual:
            \begin{itemize}[nosep]
                \item Si estado es ``Creada'': puede pasar a ``En proceso'' o ``Cancelada''.
                \item Si estado es ``En proceso'': puede pasar a ``Finalizada'' o ``Cancelada''.
            \end{itemize}
        \item El personal seleccióna el nuevo estado.
        \item El personal escribe observaciones (obligatorio para ``Cancelada'', opcional para los demas).
        \item El personal confirma el cambio.
        \item El sistema actualiza el estado de la solicitud.
        \item El sistema registra el cambio en el historial de seguimiento con: fecha, estado anterior, nuevo estado, responsable y observaciones.
        \item Si el nuevo estado es ``Finalizada'', el sistema registra la fecha de finalización.
        \item El sistema envia notificación por correo al solicitante informando del cambio de estado (CU-17).
        \item El sistema muestra mensaje de éxito.
    \end{enumerate} \\
    \hline
    \textbf{Postcondición(es):} &
    \begin{enumerate}[nosep, leftmargin=*]
        \item El estado de la solicitud esta actualizado.
        \item El historial de seguimiento tiene un nuevo registro.
        \item El solicitante fue notificado por correo.
    \end{enumerate} \\
    \hline
    \textbf{Alternativos:} &
    \textbf{FA-01: Error al enviar notificación.}
    \begin{enumerate}[nosep, leftmargin=*]
        \item[10.1] El servicio SMTP no responde.
        \item[10.2] El cambio de estado se completa éxitosamente.
        \item[10.3] La notificación queda pendiente de reintento.
    \end{enumerate}
    \vspace{0.2cm}
    \textbf{FA-02: Observaciones vacias para cancelación.}
    \begin{enumerate}[nosep, leftmargin=*]
        \item[6.1] El sistema detecta que no se proporcionaron observaciones para una cancelación.
        \item[6.2] Se muestra un mensaje de error solicitando las observaciones.
        \item[6.3] Regresa al paso 5.
    \end{enumerate} \\
    \hline
    \textbf{Req. relacionados:} & RF-05, RF-07, RF-09 \\
    \hline
\end{longtable}

% --- CU-09 ---
\subsection{CU-09: Generar documento PDF}

\begin{longtable}{|L{4cm}|L{9.5cm}|}
    \hline
    \multicolumn{2}{|c|}{\textbf{CU-09: Generar documento PDF}} \\
    \hline
    \endfirsthead
    \hline
    \endfoot

    \textbf{Fecha y Versión:} & Marzo 2026, v1.0 \\
    \hline
    \textbf{Autor:} & Equipo de desarrollo \\
    \hline
    \textbf{Descripción:} & El personal genera un documento PDF a partir de la plantilla asociada al tipo de solicitud, sustituyendo las variables definidas con los datos reales del solicitante y la solicitud. \\
    \hline
    \textbf{Actor(es) primario(s):} & Control Escolar, Responsable de Programa \\
    \hline
    \textbf{Actor(es) secundario(s):} & SIGA (datos del solicitante) \\
    \hline
    \textbf{Precondición(es):} &
    \begin{enumerate}[nosep, leftmargin=*]
        \item El usuario esta autenticado con rol de personal.
        \item La solicitud tiene un tipo de solicitud con plantilla asociada.
        \item La plantilla tiene variables definidas.
    \end{enumerate} \\
    \hline
    \textbf{Happy Path:} &
    \begin{enumerate}[nosep, leftmargin=*]
        \item El personal visualiza la solicitud (CU-07).
        \item El sistema detecta que el tipo de solicitud tiene una plantilla asociada y muestra el boton ``Generar documento''.
        \item El personal presiona ``Generar documento''.
        \item El sistema obtiene los datos del solicitante desde el SIGA (nombre, matrícula, programa, etc.).
        \item El sistema obtiene las respuestas del formulario de la solicitud.
        \item El sistema sustituye las variables de la plantilla (ej. \texttt{\{\{nombre\}\}}, \texttt{\{\{matrícula\}\}}, \texttt{\{\{fecha\}\}}) con los valores reales.
        \item El sistema genera el documento en formato PDF.
        \item El sistema ofrece la descarga del PDF al personal.
    \end{enumerate} \\
    \hline
    \textbf{Postcondición(es):} & El documento PDF generado esta disponible para descarga. \\
    \hline
    \textbf{Alternativos:} &
    \textbf{FA-01: SIGA no disponible.}
    \begin{enumerate}[nosep, leftmargin=*]
        \item[4.1] No se pueden obtener datos completos del solicitante.
        \item[4.2] El sistema muestra un mensaje de error indicando que no se puede generar el documento en este momento.
    \end{enumerate}
    \vspace{0.2cm}
    \textbf{FA-02: Variable sin valor.}
    \begin{enumerate}[nosep, leftmargin=*]
        \item[6.1] Una variable de la plantilla no tiene un valor correspondiente.
        \item[6.2] El sistema deja la variable vacia o muestra un marcador ``[SIN DATO]''.
        \item[6.3] Continua con la generación del PDF.
    \end{enumerate} \\
    \hline
    \textbf{Req. relacionados:} & RF-03, RF-09, RNF-02 \\
    \hline
\end{longtable}

% --- CU-10 ---
\subsection{CU-10: Gestiónar tipos de solicitud}

\begin{longtable}{|L{4cm}|L{9.5cm}|}
    \hline
    \multicolumn{2}{|c|}{\textbf{CU-10: Gestiónar tipos de solicitud}} \\
    \hline
    \endfirsthead
    \hline
    \endfoot

    \textbf{Fecha y Versión:} & Marzo 2026, v1.0 \\
    \hline
    \textbf{Autor:} & Equipo de desarrollo \\
    \hline
    \textbf{Descripción:} & El administrador gestióna el catálogo de tipos de solicitud, pudiendo crear nuevos tipos, editar los existentes o eliminarlos. Cada tipo define su nombre, descripción, rol responsable, si requiere pago y los roles que pueden crearlo. \\
    \hline
    \textbf{Actor(es) primario(s):} & Administrador \\
    \hline
    \textbf{Actor(es) secundario(s):} & --- \\
    \hline
    \textbf{Precondición(es):} &
    \begin{enumerate}[nosep, leftmargin=*]
        \item El usuario esta autenticado con rol de administrador.
    \end{enumerate} \\
    \hline
    \textbf{Happy Path (Crear):} &
    \begin{enumerate}[nosep, leftmargin=*]
        \item El administrador accede a la sección ``Catálogo de solicitudes''.
        \item El sistema muestra la lista de tipos de solicitud existentes.
        \item El administrador presiona ``Agregar tipo''.
        \item El sistema presenta un formulario con los campos:
            \begin{itemize}[nosep]
                \item Nombre del tipo de solicitud.
                \item Descripción.
                \item Rol responsable de atenderla (selección).
                \item Requiere pago (si/no).
                \item Roles que pueden crear este tipo (alumno, docente o ambos).
            \end{itemize}
        \item El administrador llena los campos y presiona ``Guardar''.
        \item El sistema valida los datos y crea el tipo de solicitud.
        \item El sistema muestra mensaje de éxito y regresa a la lista.
    \end{enumerate} \\
    \hline
    \textbf{Postcondición(es):} & El tipo de solicitud esta disponible en el catálogo para los roles especificados. \\
    \hline
    \textbf{Alternativos:} &
    \textbf{FA-01: Editar tipo existente.}
    \begin{enumerate}[nosep, leftmargin=*]
        \item[3.1] El administrador seleccióna ``Editar'' en un tipo existente.
        \item[3.2] El sistema carga el formulario con los datos actuales.
        \item[3.3] El administrador modifica los campos deseados y guarda.
    \end{enumerate}
    \vspace{0.2cm}
    \textbf{FA-02: Eliminar tipo.}
    \begin{enumerate}[nosep, leftmargin=*]
        \item[3.1] El administrador seleccióna ``Eliminar'' en un tipo.
        \item[3.2] El sistema verifica si existen solicitudes asociadas a ese tipo.
        \item[3.3a] Si no hay solicitudes: se elimina el tipo.
        \item[3.3b] Si hay solicitudes: se muestra un mensaje indicando que no se puede eliminar porque tiene solicitudes asociadas.
    \end{enumerate}
    \vspace{0.2cm}
    \textbf{FA-03: Nombre duplicado.}
    \begin{enumerate}[nosep, leftmargin=*]
        \item[6.1] Ya existe un tipo con el mismo nombre.
        \item[6.2] El sistema muestra un error de duplicidad.
    \end{enumerate} \\
    \hline
    \textbf{Req. relacionados:} & RF-01 \\
    \hline
\end{longtable}

% --- CU-11 ---
\subsection{CU-11: Gestiónar formularios dinámicos}

\begin{longtable}{|L{4cm}|L{9.5cm}|}
    \hline
    \multicolumn{2}{|c|}{\textbf{CU-11: Gestiónar formularios dinámicos}} \\
    \hline
    \endfirsthead
    \hline
    \endfoot

    \textbf{Fecha y Versión:} & Marzo 2026, v1.0 \\
    \hline
    \textbf{Autor:} & Equipo de desarrollo \\
    \hline
    \textbf{Descripción:} & El administrador configura los campos del formulario asociado a un tipo de solicitud. Puede agregar, editar, reordenar y eliminar campos, definiendo su tipo, etiqueta, si es requerido y opciones adicionales. \\
    \hline
    \textbf{Actor(es) primario(s):} & Administrador \\
    \hline
    \textbf{Actor(es) secundario(s):} & --- \\
    \hline
    \textbf{Precondición(es):} &
    \begin{enumerate}[nosep, leftmargin=*]
        \item El usuario esta autenticado con rol de administrador.
        \item Existe al menos un tipo de solicitud en el catálogo.
    \end{enumerate} \\
    \hline
    \textbf{Happy Path:} &
    \begin{enumerate}[nosep, leftmargin=*]
        \item El administrador accede a la sección ``Formularios'' y seleccióna un tipo de solicitud.
        \item El sistema muestra los campos actuales del formulario (si existen), ordenados por su campo ``orden''.
        \item El administrador presiona ``Agregar campo''.
        \item El sistema presenta un formulario con:
            \begin{itemize}[nosep]
                \item Etiqueta del campo (lo que vera el usuario).
                \item Tipo de campo: texto, area de texto, número, fecha, selección (dropdown), archivo.
                \item Requerido (si/no).
                \item Opciones (solo para tipo ``selección'', separadas por coma).
                \item Orden de aparicion (número).
            \end{itemize}
        \item El administrador llena la configuración y presiona ``Guardar''.
        \item El sistema valida que no haya ordenes duplicados y crea el campo.
        \item El sistema actualiza la vista del formulario con el nuevo campo.
    \end{enumerate} \\
    \hline
    \textbf{Postcondición(es):} & El formulario queda actualizado y será presentado a los solicitantes al crear una solicitud de ese tipo. \\
    \hline
    \textbf{Alternativos:} &
    \textbf{FA-01: Editar campo existente.}
    \begin{enumerate}[nosep, leftmargin=*]
        \item[3.1] El administrador seleccióna ``Editar'' en un campo.
        \item[3.2] Modifica la configuración y guarda.
    \end{enumerate}
    \vspace{0.2cm}
    \textbf{FA-02: Eliminar campo.}
    \begin{enumerate}[nosep, leftmargin=*]
        \item[3.1] El administrador seleccióna ``Eliminar'' en un campo.
        \item[3.2] El sistema solicita confirmación.
        \item[3.3] Se elimina el campo del formulario.
    \end{enumerate}
    \vspace{0.2cm}
    \textbf{FA-03: Orden duplicado.}
    \begin{enumerate}[nosep, leftmargin=*]
        \item[6.1] Ya existe un campo con el mismo número de orden.
        \item[6.2] El sistema muestra un error y solicita un orden diferente.
    \end{enumerate} \\
    \hline
    \textbf{Req. relacionados:} & RF-02 \\
    \hline
\end{longtable}

% --- CU-12 ---
\subsection{CU-12: Gestiónar plantillas}

\begin{longtable}{|L{4cm}|L{9.5cm}|}
    \hline
    \multicolumn{2}{|c|}{\textbf{CU-12: Gestiónar plantillas de documentos}} \\
    \hline
    \endfirsthead
    \hline
    \endfoot

    \textbf{Fecha y Versión:} & Marzo 2026, v1.0 \\
    \hline
    \textbf{Autor:} & Equipo de desarrollo \\
    \hline
    \textbf{Descripción:} & El administrador crea, edita o elimina plantillas de documentos que se utilizan para generar PDFs. Las plantillas contienen variables sustituibles como \texttt{\{\{nombre\}\}}, \texttt{\{\{matrícula\}\}}, \texttt{\{\{programa\}\}}, \texttt{\{\{fecha\}\}}. \\
    \hline
    \textbf{Actor(es) primario(s):} & Administrador \\
    \hline
    \textbf{Actor(es) secundario(s):} & --- \\
    \hline
    \textbf{Precondición(es):} &
    \begin{enumerate}[nosep, leftmargin=*]
        \item El usuario esta autenticado con rol de administrador.
        \item Existe al menos un tipo de solicitud que requiera plantilla.
    \end{enumerate} \\
    \hline
    \textbf{Happy Path:} &
    \begin{enumerate}[nosep, leftmargin=*]
        \item El administrador accede a la sección ``Plantillas''.
        \item El sistema muestra las plantillas existentes con su tipo de solicitud asociado.
        \item El administrador presiona ``Crear plantilla''.
        \item El sistema presenta un formulario con:
            \begin{itemize}[nosep]
                \item Tipo de solicitud asociado (selección).
                \item Nombre de la plantilla.
                \item Contenido de la plantilla con soporte para variables (ej. \texttt{\{\{nombre\}\}}, \texttt{\{\{matrícula\}\}}).
                \item Lista de variables disponibles (referencia).
            \end{itemize}
        \item El administrador configura la plantilla y presiona ``Guardar''.
        \item El sistema valida y almacena la plantilla.
        \item El sistema muestra mensaje de éxito.
    \end{enumerate} \\
    \hline
    \textbf{Postcondición(es):} & La plantilla queda disponible para la generación de documentos PDF en solicitudes del tipo asociado. \\
    \hline
    \textbf{Alternativos:} &
    \textbf{FA-01: Editar plantilla existente.}
    \begin{enumerate}[nosep, leftmargin=*]
        \item[3.1] El administrador seleccióna ``Editar'' en una plantilla.
        \item[3.2] Modifica el contenido y guarda.
    \end{enumerate}
    \vspace{0.2cm}
    \textbf{FA-02: Eliminar plantilla.}
    \begin{enumerate}[nosep, leftmargin=*]
        \item[3.1] El administrador seleccióna ``Eliminar''.
        \item[3.2] Se solicita confirmación y se elimina.
    \end{enumerate} \\
    \hline
    \textbf{Req. relacionados:} & RF-03 \\
    \hline
\end{longtable}

% --- CU-13 ---
\subsection{CU-13: Gestiónar lista de mentores}

\begin{longtable}{|L{4cm}|L{9.5cm}|}
    \hline
    \multicolumn{2}{|c|}{\textbf{CU-13: Gestiónar lista de mentores}} \\
    \hline
    \endfirsthead
    \hline
    \endfoot

    \textbf{Fecha y Versión:} & Marzo 2026, v1.0 \\
    \hline
    \textbf{Autor:} & Equipo de desarrollo \\
    \hline
    \textbf{Descripción:} & El administrador gestióna la lista de alumnos mentores activos. Los mentores estan exentos de adjuntar comprobante de pago en solicitudes que lo requieran. La lista se puede gestiónar de forma manual o mediante carga de CSV. \\
    \hline
    \textbf{Actor(es) primario(s):} & Administrador \\
    \hline
    \textbf{Actor(es) secundario(s):} & --- \\
    \hline
    \textbf{Precondición(es):} &
    \begin{enumerate}[nosep, leftmargin=*]
        \item El usuario esta autenticado con rol de administrador.
    \end{enumerate} \\
    \hline
    \textbf{Happy Path (Manual):} &
    \begin{enumerate}[nosep, leftmargin=*]
        \item El administrador accede a la sección ``Gestión de mentores''.
        \item El sistema muestra la lista actual de alumnos mentores con: matrícula, nombre (consultado al SIGA si esta disponible) y fecha de alta.
        \item El administrador presiona ``Agregar mentor''.
        \item El administrador ingresa la matrícula del alumno.
        \item El sistema valida que la matrícula exista (consulta al SIGA si es posible).
        \item El sistema agrega al alumno a la lista de mentores.
        \item El sistema muestra mensaje de éxito.
    \end{enumerate} \\
    \hline
    \textbf{Postcondición(es):} & El alumno queda registrado como mentor activo y será exento de pago en solicitudes que lo requieran. \\
    \hline
    \textbf{Alternativos:} &
    \textbf{FA-01: Carga masiva por CSV.}
    \begin{enumerate}[nosep, leftmargin=*]
        \item[3.1] El administrador presiona ``Cargar CSV''.
        \item[3.2] El administrador seleccióna un archivo CSV con una columna de matrículas.
        \item[3.3] El sistema lee el archivo y procesa cada matrícula.
        \item[3.4] El sistema muestra un resumen: cuantos se agregaron, cuantos ya existian, cuantos tuvieron error.
    \end{enumerate}
    \vspace{0.2cm}
    \textbf{FA-02: Eliminar mentor.}
    \begin{enumerate}[nosep, leftmargin=*]
        \item[2.1] El administrador seleccióna ``Eliminar'' junto a un mentor.
        \item[2.2] Se solicita confirmación.
        \item[2.3] Se elimina al alumno de la lista de mentores.
    \end{enumerate}
    \vspace{0.2cm}
    \textbf{FA-03: Matricula duplicada.}
    \begin{enumerate}[nosep, leftmargin=*]
        \item[6.1] La matrícula ya esta en la lista de mentores.
        \item[6.2] El sistema muestra un mensaje indicando que ya esta registrado.
    \end{enumerate}
    \vspace{0.2cm}
    \textbf{FA-04: Matricula no encontrada.}
    \begin{enumerate}[nosep, leftmargin=*]
        \item[5.1] La matrícula no existe en el SIGA.
        \item[5.2] El sistema muestra una advertencia pero permite agregarlo (la validación con SIGA es complementaria).
    \end{enumerate} \\
    \hline
    \textbf{Req. relacionados:} & RF-11 \\
    \hline
\end{longtable}

% --- CU-14 ---
\subsection{CU-14: Ver dashboard y métricas}

\begin{longtable}{|L{4cm}|L{9.5cm}|}
    \hline
    \multicolumn{2}{|c|}{\textbf{CU-14: Ver dashboard y métricas}} \\
    \hline
    \endfirsthead
    \hline
    \endfoot

    \textbf{Fecha y Versión:} & Marzo 2026, v1.0 \\
    \hline
    \textbf{Autor:} & Equipo de desarrollo \\
    \hline
    \textbf{Descripción:} & El administrador visualiza un panel con estadísticas generales del sistema: total de solicitudes, distribucion por tipo, por estado y por período de tiempo. \\
    \hline
    \textbf{Actor(es) primario(s):} & Administrador \\
    \hline
    \textbf{Actor(es) secundario(s):} & --- \\
    \hline
    \textbf{Precondición(es):} &
    \begin{enumerate}[nosep, leftmargin=*]
        \item El usuario esta autenticado con rol de administrador.
        \item Existe al menos una solicitud en el sistema.
    \end{enumerate} \\
    \hline
    \textbf{Happy Path:} &
    \begin{enumerate}[nosep, leftmargin=*]
        \item El administrador accede a la sección ``Dashboard''.
        \item El sistema calcula y muestra:
            \begin{itemize}[nosep]
                \item Total de solicitudes.
                \item Solicitudes por estado (gráficos).
                \item Solicitudes por tipo (gráficos).
                \item Solicitudes del dia, la semana y el mes.
                \item Solicitudes por area/responsable.
            \end{itemize}
        \item El administrador puede filtrar por rango de fechas.
        \item El administrador puede selecciónar exportar los datos (CU-15).
    \end{enumerate} \\
    \hline
    \textbf{Postcondición(es):} & El administrador visualiza las métricas actualizadas. \\
    \hline
    \textbf{Req. relacionados:} & RF-08, RNF-05 \\
    \hline
\end{longtable}

% --- CU-15 ---
\subsection{CU-15: Exportar reportes}

\begin{longtable}{|L{4cm}|L{9.5cm}|}
    \hline
    \multicolumn{2}{|c|}{\textbf{CU-15: Exportar reportes}} \\
    \hline
    \endfirsthead
    \hline
    \endfoot

    \textbf{Fecha y Versión:} & Marzo 2026, v1.0 \\
    \hline
    \textbf{Autor:} & Equipo de desarrollo \\
    \hline
    \textbf{Descripción:} & El administrador exporta las métricas y datos de solicitudes en formato CSV o PDF para su análisis o presentación fuera del sistema. \\
    \hline
    \textbf{Actor(es) primario(s):} & Administrador \\
    \hline
    \textbf{Actor(es) secundario(s):} & --- \\
    \hline
    \textbf{Precondición(es):} &
    \begin{enumerate}[nosep, leftmargin=*]
        \item El usuario esta autenticado con rol de administrador.
        \item Existen datos para exportar.
    \end{enumerate} \\
    \hline
    \textbf{Happy Path:} &
    \begin{enumerate}[nosep, leftmargin=*]
        \item El administrador accede al dashboard (CU-14) o a la sección de reportes.
        \item El administrador seleccióna el formato de exportación: CSV o PDF.
        \item El administrador opcionalmente seleccióna filtros (rango de fechas, tipo, estado).
        \item El administrador presiona ``Exportar''.
        \item El sistema genera el archivo en el formato selecciónado.
        \item El sistema ofrece la descarga del archivo.
    \end{enumerate} \\
    \hline
    \textbf{Postcondición(es):} & El archivo de reporte se descarga éxitosamente. \\
    \hline
    \textbf{Req. relacionados:} & RNF-05 \\
    \hline
\end{longtable}

% --- CU-17 ---
\subsection{CU-17: Enviar notificación por correo}

\begin{longtable}{|L{4cm}|L{9.5cm}|}
    \hline
    \multicolumn{2}{|c|}{\textbf{CU-17: Enviar notificación por correo}} \\
    \hline
    \endfirsthead
    \hline
    \endfoot

    \textbf{Fecha y Versión:} & Marzo 2026, v1.0 \\
    \hline
    \textbf{Autor:} & Equipo de desarrollo \\
    \hline
    \textbf{Descripción:} & El sistema envia automáticamente notificaciónes por correo electrónico en eventos relevantes del ciclo de vida de una solicitud. No es un caso de uso iniciado por un actor humano; es disparado por otros casos de uso. \\
    \hline
    \textbf{Actor(es) primario(s):} & Sistema (automático) \\
    \hline
    \textbf{Actor(es) secundario(s):} & Servicio SMTP \\
    \hline
    \textbf{Precondición(es):} &
    \begin{enumerate}[nosep, leftmargin=*]
        \item Se ha producido un evento que requiere notificación (creación de solicitud o cambio de estado).
        \item El destinatario tiene un correo electrónico registrado.
    \end{enumerate} \\
    \hline
    \textbf{Happy Path:} &
    \begin{enumerate}[nosep, leftmargin=*]
        \item El sistema detecta un evento de notificación:
            \begin{itemize}[nosep]
                \item Solicitud creada: notificar al personal responsable.
                \item Estado cambiado: notificar al solicitante.
            \end{itemize}
        \item El sistema construye el correo con:
            \begin{itemize}[nosep]
                \item Destinatario(s).
                \item Asunto descriptivo (ej. ``Nueva solicitud: [tipo] - Folio [folio]'').
                \item Cuerpo con información relevante (folio, tipo, estado, enlace al sistema).
            \end{itemize}
        \item El sistema envia el correo a través del servicio SMTP.
        \item El envio se registra internamente.
    \end{enumerate} \\
    \hline
    \textbf{Postcondición(es):} & El correo fue enviado éxitosamente al destinatario. \\
    \hline
    \textbf{Alternativos:} &
    \textbf{FA-01: Servicio SMTP no disponible.}
    \begin{enumerate}[nosep, leftmargin=*]
        \item[3.1] La conexión al SMTP falla.
        \item[3.2] El sistema registra el error.
        \item[3.3] La notificación queda marcada como pendiente para reintento posterior.
        \item[3.4] El flujo principal del caso de uso que disparo la notificación no se ve afectado.
    \end{enumerate} \\
    \hline
    \textbf{Req. relacionados:} & RF-07 \\
    \hline
\end{longtable}
